\section{Scriptural Date and Location}
\label{sec:date-location}
{\small Each passage the Mass reads or sings from Scripture, once, in canonical
order; beneath each, its setting, attribution and dating evidence.\par}
\medskip

\begin{dossiertable}
Alleluia & Ps 107:2 (Heb.\ 108) & David, by the title; the psalm repeats the
closes of Pss~56 and 59 &
\chronodate{alleluia}{\chronologyannotation{alleluia}}\\
\cmidrule(lr){1-4}
\dossierprose{Titled in the Vulgate \latin{Canticum Psalmi ipsi David}, ``A
canticle of a psalm for David himself'', which names no occasion. Its verses
repeat Ps~56:8--12 and 59:7--14, whose titles set them in David's flight from
Saul into the cave and in the war in which Joab slew of Edom in the vale of the
saltpits; those occasions belong to the two source psalms, and the chronology
record holds none for Ps~107. The reign shown with the attribution is David's
in the ``usual chronology'' that the \work{Catholic Encyclopedia}'s ``David''
(vol.~4, 1908) reports, adding that recent writers date it thirty to fifty
years later: it is the attribution's reference point, not the psalm's date. The
composition date is the bound that the introduction to the Psalms in the
\work{New American Bible, Revised Edition} sets for the whole Psalter, within
which, it says, no individual psalm can be dated securely.}
\midrule
Communion & Ps 118:49--50 (Heb.\ 119) & The psalmist, ``a sojourner on the
earth'' (v.~19), praying in the strophe \latin{Zain} &
\chronodate{communion}{\chronologyannotation{communion}}\\
\cmidrule(lr){1-4}
\dossierprose{Titled \latin{Alleluia} in the Vulgate, with no author named. The
Davidic attribution shown is the tradition's for the whole psalm, marked
disputed, and its reign is again the attribution's reference point; the record
holds no occasion, and the composition date is the Psalter's bound. The
Introit's verse, Ps~118:1, stands in the same psalm and is dated with the
Introit below.}
\midrule
Offertory & Ps 136:1 (Heb.\ 137) & The captives of Juda by the rivers of
Babylon, on the Euphrates, after the fall of Jerusalem; the Vulgate's title
names David and Jeremias &
\chronodate{offertory}{\chronologyannotation{offertory}}\\
\cmidrule(lr){1-4}
\dossierprose{Titled \latin{Psalmus David, Ieremiae}, a title the Missal leaves
out; St Hilary notes that among the Hebrew elders the psalm has none. The record
gives it a setting rather than an author: the lament of the Babylonian captives,
``For the time of Jeremias, and the captivity of Babylon'' (Haydock), and the
seventy years of servitude to Babylon (Jer 25:11--12), which are a length of
time and not a date. The dates under ``Historical background'' belong to the
destruction of Jerusalem under Sedecias, which the lament follows, and not to
the psalm or its writing; the year given in the reckoning from the creation of
the world is Haydock's report of Ussher, left unconverted. The composition date
is the Psalter's bound.}
\midrule
Gradual & Ps 144:15--16 (Heb.\ 145) & David, by the title; a hymn of God's
kingship and of his care for every creature &
\chronodate{gradual}{\chronologyannotation{gradual}}\\
\cmidrule(lr){1-4}
\dossierprose{Titled \latin{Laudatio ipsi David}, ``Praise, for David himself'',
without an occasion. The reign shown is the attribution's reference point, not
the psalm's date, and the composition date is the Psalter's bound.}
\midrule
Introit & Dan 3:31, 29 and 35, as printed; the verse Ps 118:1 (Heb.\ 119) &
The furnace in the province of Babylon, under Nabuchodonosor (Dan 3:1, 20);
Azarias prays for the exiled people &
\chronodate{introit}{\chronologyannotation{introit}}\\
\cmidrule(lr){1-4}
\dossierevent{Narrated event: the three young men, cast into the burning furnace
for refusing to adore the king's golden statue set up in the plain of Dura,
walk in the flame praising God, and Azarias prays (Dan 3:1--25); the narrative
gives the scene no year.}
\cmidrule(lr){1-4}
\dossierprose{The antiphon's Daniel and its psalm verse are dated differently,
so no single date stands in the cell. For the prayer, the record holds a
judgment of composition without a year: the \work{Catholic Encyclopedia}'s
``Book of Daniel'' (vol.~4, 1908) ascribes the Prayer of Azarias and the Song of
the Three Children (Dan 3:24--90) to an unknown Jewish author, ``most likely
long after the Exile'', and the record sets the prayer's world within the
seventy years of servitude to Babylon, a length of time and not a date. St
Jerome reports that these verses are not in the Hebrew and that the Churches
read Daniel in Theodotion's version. For the verse, the record gives the
Psalter's composition bound, Ps~118:1: before c.~165~B.C.; and, as the
reference point of the disputed traditional attribution to David, his reign in
the usual chronology, Ps~118:1: 1055--1015~B.C.}
\midrule
Gospel & Jn 4:46--53 & Writing: St John the Apostle, by the Missal's heading.
Event: the word at Cana, the cure at Capharnaum, in Galilee (4:46) &
\chronodate{gospel}{\chronologyannotation{gospel}}\\
\cmidrule(lr){1-4}
\dossierevent{Narrated event: the healing of the ruler's son, ``the second
miracle that Jesus did, when he was come out of Judea into Galilee'' (4:54),
after two days among the Samaritans (4:40--43); the word is spoken at Cana and
the fever leaves the boy at Capharnaum at the same hour, the seventh hour of
the day before the servants' report (4:52).}
\cmidrule(lr){1-4}
\dossierprose{The event's year is derived. The \work{Catholic Encyclopedia}'s
``Jesus Christ'' (vol.~8, 1910), by A.~J. Maas, places the return to Cana in a
period it reckons from the Passover of one year to about Pentecost of the next,
counted from the founding of Rome, and the displayed years convert that
reckoning. Traditional attribution: St John the Apostle, whom the Missal's
heading names. The two composition dates shown are the encyclopedia's, both marked
disputed: the first, approximate, from ``The New Testament'' (vol.~14, 1912),
and the second from ``The Gospel of Saint John'' (vol.~8, 1910).}
\midrule
% booktabs leaves a legal page break after each partial rule; the Epistle's
% dossier is kept whole by starting it where the table already stands.
Epistle & Eph 5:15--21 & Paul, ``a prisoner in the Lord'' (4:1), in Caesarea or
Rome; to ``the saints who are at Ephesus'' (1:1) &
\chronodate{epistle}{\chronologyannotation{epistle}}\\
\cmidrule(lr){1-4}
\dossierprose{Traditional attribution: St Paul, as the Missal's heading names
him, who names himself (1:1) and writes as a prisoner (3:1; 4:1) and ``an
ambassador in a chain'' (6:20). The first range shown is that of the
\work{Catholic Encyclopedia}'s ``Epistle to the Ephesians'' (vol.~5, 1909),
where of the letters of the captivity ``whether they were produced in Caesarea
or in Rome \dots\ is still a much mooted question''. The year is from the table
in its ``St.\ Paul'' (vol.~11, 1911), which sets the letter in the Roman
captivity with Philemon, Colossians and Philippians; both are marked disputed.
The comparison shown last is the modern critical horizon: the \work{New American
Bible, Revised Edition}'s introduction to Ephesians offers, as one possible
explanation of the letter's style and teaching, a later disciple developing
Paul's thought, around A.D.~80--100. That is a conditional hypothesis and not
the default answer.}
\end{dossiertable}
